\documentclass[11pt]{article}

\usepackage[margin=0.82in]{geometry}
\usepackage[T1]{fontenc}
\usepackage[utf8]{inputenc}
\usepackage{lmodern}
\usepackage{amsmath,amssymb}
\usepackage{booktabs}
\usepackage{array}
\usepackage{graphicx}
\usepackage{subcaption}
\usepackage{placeins}
\usepackage{float}
\floatstyle{ruled}
\newfloat{algorithm}{htbp}{loa}
\floatname{algorithm}{Algorithm}

\usepackage{xcolor}
\usepackage{microtype}
\usepackage{hyperref}
\setlength{\parindent}{0pt}
\setlength{\parskip}{0.55em}

\hypersetup{colorlinks=true,linkcolor=blue!55!black,urlcolor=blue!55!black,citecolor=blue!55!black}

\IfFileExists{figures/mpl-surface3d-liquidity-only-low-z-v3.pdf}
  {\newcommand{\FigDir}{figures}}
  {\IfFileExists{../figures/mpl-surface3d-liquidity-only-low-z-v3.pdf}
    {\newcommand{\FigDir}{../figures}}{\newcommand{\FigDir}{../../figures}}}
\newcommand{\Aeff}{\mathcal A_{\mathrm{eff}}}
\newcommand{\Liq}{\mathcal L}
\newcommand{\Fresp}{\mathcal F}
\newcommand{\sigloc}{\sigma_{\mathrm{loc}}}

\title{Computational Supplement:\\Reaction-Boundary Volatility Surface Simulations}
\author{Chris Angstmann \and Tim Gebbie}
\date{Supplement v1.4.0 --- Independent response and forcing exponents\\22 November 2026}

\begin{document}
\maketitle

\begin{abstract}
The frozen reaction boundary and the translating liquidity profile give different
responses to persistent order flow. We examine how finite resolution,
resilience and liquidity enter the structural volatility surfaces, and then
calculate the response of an aligned translating root to prescribed liquidity
reassessment. A finite lower rate cutoff leaves a positive zero-frequency
mobility. The accuracy with which fast rates can be eliminated is measured
over an explicitly retained band of operational frequencies. We then measure forcing memory and integrated root response independently over a shared operational-lag window and compare the matched and offset rate laws.
\end{abstract}
\tableofcontents

\section{Scientific scope and background}

This supplementary note describes a simulation study of volatility surfaces
generated from the reaction-boundary closure. The purpose is not calibration.
Instead, the figures show how structural order-book channels can produce
stock-price and maturity dependence in a projected local-volatility surface.
The simulations are intended as a visual bridge between the theoretical
closure and the kinds of surfaces one would later compare with calibrated
local-volatility data.

The reference for the frozen response is Angstmann and Gebbie,
\href{https://arxiv.org/abs/2607.05011v4}{arXiv:2607.05011v4}.
The translating field and distributed reassessment below are auxiliary
extensions. Throughout the response calculations \(u\) is operational
time. In the surface plots \(m\) denotes log-moneyness. We use \(m_u\)
for the signed input to the moving boundary and \(r\) for a reassessment
rate in inverse operational time.

\subsection{Theory-to-computation map}
\begin{center}
\small
\begin{tabular}{p{.23\linewidth}p{.31\linewidth}p{.36\linewidth}}
\toprule
Object & Observable & Independent comparison \\
\midrule
Frozen imbalance zero & Increment variance at fixed cutoff & Spectral and time-domain quadratures at identical parameters \\
Translating profile & Root, neutral eigenmode and shape band & Lattice pole, observed zero and weak field residual \\
Prescribed reassessment & Velocity transfer and integrated root response & Truncated integral, mode evolution and continuum relaxation \\
Fast reassessment rates & Response on the retained frequency band & Full memory and successive local approximations \\
Response balance & Input spectrum and integrated root response & Independent path estimates, inversion refinements and offset controls \\
\bottomrule
\end{tabular}
\end{center}

\section{Reaction-Boundary Volatility Closure}

The operational variance closure used in the simulation is
\begin{equation}
 a_u^{(\Delta)}(S,u)
 \simeq
 \left[
 \frac{\Aeff(S,u)}
      {\Liq(S,u)^2}
 \right]
 \Delta^{-\gamma(S,u)}
 \Fresp_{\gamma(S,u)}\!\left(\nu(S,u)\Delta\right).
 \label{eq:operational-closure}
\end{equation}
Here \(S\) denotes the stock-price coordinate, \(u\) is operational time,
\(\Aeff\) is an effective signed-forcing scale, \(\Liq\) is the local
liquidity slope, \(\Delta\) is the operational coarse-graining scale,
\(\gamma\) is the long-memory exponent, and
\(z=\nu(S,u)\Delta\) is the dimensionless resilience-response argument.

Under a deterministic activity clock, the projected local variance is written
as
\begin{equation}
 \sigloc^2(S,t)
 =
 \alpha(t)\,
 \Xi(S,U(t);\Delta),
 \label{eq:projected-local-variance}
\end{equation}
where \(\alpha(t)=U'(t)\) maps operational-time variance into calendar time. In the
simulations, \(\Xi\) is evaluated using the right-hand side of
Eq.~\eqref{eq:operational-closure}. The plotted quantity is projected local
volatility,
\begin{equation}
 \sigloc(m,T)=\sqrt{\sigloc^2(m,T)},
 \label{eq:plotted-volatility}
\end{equation}
on a grid in log-moneyness \(m=\ln(S/S_0)\) and maturity \(T\).

For a stationary frozen root \(Y\), write
\(V_\Delta=\mathrm{Var}[Y(u+\Delta)-Y(u)]\) and \(V_Y=\mathrm{Var}[Y(u)]\).
The unconditional variance rate is \(V_\Delta/\Delta\). In the normalisation
of the frozen field, \(\Aeff=\mathcal A_m/D\), where \(D\) is liquidity
diffusivity and \(\mathcal A_m\) sets the forcing covariance. Under Gaussian
stationarity the conditional increment second moment is
\[
 V_\Delta+\frac{V_\Delta^2}{4V_Y}\left(\frac{y^2}{V_Y}-1\right)
 \quad\text{at }Y(u)=y .
\]
The conditional increment mean is \(-y V_\Delta/(2V_Y)\) and the
conditional variance is \(V_\Delta-V_\Delta^2/(4V_Y)\) for \(V_Y>0\).
Adding the squared conditional mean gives the displayed second moment.
This follows by Gaussian regression using
\(\mathrm{Cov}(Y(u+\Delta)-Y(u),Y(u))=-V_\Delta/2\).
When \(V_\Delta\to2V_Y\), the conditional moment tends to \(V_Y+y^2\).
A pricing interpretation of Eq.~\eqref{eq:projected-local-variance}
requires state conditioning and a choice of operational parameters under
the pricing measure.

\section{Response-Regime Approximations}

For the regularised causal kernel
\(g^{(\tau_0)}(h)=e^{-\nu(h+\tau_0)}/\sqrt{4\pi D(h+\tau_0)}\)
on \(h\ge0\), take \(C_m(h)=\mathcal A_m(|h|+\tau_0)^{-\gamma}\).
At fixed \(\tau_0>0\) the stationary filter is mean-square differentiable
and \(V_h=O(h^2)\) as \(h\downarrow0\). The low-\(z\) scaling interval
requires \(\tau_0\ll\Delta\ll\nu^{-1}\). The zero-cutoff high-\(z\)
prefactor also requires \(\nu\tau_0\ll1\). In 48 finite-cutoff cases,
direct covariance and spectral quadratures gave a maximum relative
discrepancy of \(3.84\times10^{-11}\).

Both asymptotic response regimes are simulated. The low-\(z\) branch uses
\begin{equation}
 \Fresp_\gamma(z)
 \sim
 \frac{\tan(\pi\gamma/2)}{2(1-\gamma)},
 \qquad z\ll1,
 \label{eq:low-z-response}
\end{equation}
which removes the leading-order dependence on \(z\) when \(\gamma\) is held
fixed. The high-\(z\) branch uses
\begin{equation}
 \Fresp_\gamma(z)
 \sim
 K_\gamma^{(0)}z^{\gamma-1},
 \qquad z\gg1,
 \label{eq:high-z-response}
\end{equation}
where
\begin{equation}
 K_\gamma^{(0)}
 =
 \frac{
 \Gamma(1-\gamma)\sin(\pi\gamma/2)
 \Gamma(\gamma/2)\Gamma((1-\gamma)/2)
 }{2\pi^{3/2}}.
 \label{eq:k-gamma}
\end{equation}
Because \(0<\gamma<1\), the exponent \(\gamma-1\) is negative in the
high-\(z\) branch. Smaller \(z\), corresponding here to weaker or slower
resilience at fixed \(\Delta\), therefore raises the response multiplier in
that branch.

\begin{table}[ht]
\centering
\small
\caption{Response-function regimes used in the simulations.}
\label{tab:response-regimes}
\begin{tabular}{p{0.19\linewidth}p{0.23\linewidth}p{0.37\linewidth}p{0.11\linewidth}}
\toprule
Regime & Condition & Approximation & Depends on \(z\) \\
\midrule
Low \(z\) & \(z=\nu\Delta\ll1\) &
\(\tan(\pi\gamma/2)/(2(1-\gamma))\) & No \\
High \(z\) & \(z=\nu\Delta\gg1\) &
\(K_\gamma^{(0)}z^{\gamma-1}\) & Yes \\
\bottomrule
\end{tabular}
\end{table}

\subsection{Independent evaluations at finite cutoff}
With two-sided angular-frequency convention
\(C_m(h)=(2\pi)^{-1}\int_{\mathbb R}e^{i\omega h}S_m(\omega)\,d\omega\),
the causal frozen filter gives
\[
 V_\Delta=\frac{1}{2\pi\Liq^2}\int_{\mathbb R}
 2[1-\cos(\omega\Delta)]\,|\widehat g^{(\tau_0)}(\omega)|^2
 S_m(\omega)\,d\omega .
\]
For \(s=\nu+i\omega\), its Fourier transform is
\(e^{-\nu\tau_0}\operatorname{erfcx}(\sqrt{s\tau_0})/(2\sqrt{Ds})\),
using the square root with positive real part. Independently define
\(W(h)=\int_0^\infty g^{(\tau_0)}(r)g^{(\tau_0)}(r+h)\,dr\) for \(h\ge0\).
Integrating the increment covariance directly gives
\[
 V_\Delta=\frac{2}{\Liq^2}\int_0^\infty W(h)
 [2C_m(h)-C_m(|h-\Delta|)-C_m(h+\Delta)]\,dh .
\]
At fixed cutoff it is numerically important to resolve both the lag cusp
and the cancellation of the symmetric covariance difference. On the
\(\gamma=1/2\) comparisons, the measured variance slopes were 1.999812
for \(10^{-5}\le\Delta\le10^{-4}\), \(\tau_0=0.1\); 0.500374 for
\(10^{-4}\le\Delta\le10^{-2}\), \(\tau_0=10^{-12}\); and 0.001321 for
\(10^4\le\Delta\le10^6\), \(\tau_0=0.1\), all at unit resilience.
\begin{algorithm}[htbp]
\caption{Stationary frozen increment variance at finite resolution}
\small\textbf{Input:} \(\gamma,\nu,\tau_0,\Delta,D,\Liq,\mathcal A_m\), quadrature tolerance.\\
\textbf{Output:} Independent values of \(V_\Delta\), the level variance and refinement error.
\begin{enumerate}\setlength{\itemsep}{0.1em}
\item Evaluate the spectral transfer with the scaled complementary error function. Resolve the increment filter near the origin using \(1-\cos x=2\sin^2(x/2)\), and integrate the oscillatory tail separately.
\item In time, integrate the kernel autocorrelation. Split at \(h=\Delta\), resolve the \(\tau_0\)-scale cusp and use a symmetric expansion where the covariance difference loses precision.
\item Integrate the level covariance separately, tighten both tolerances, and compare the two increment estimates before using the conditional-moment formula.
\end{enumerate}
\end{algorithm}
\FloatBarrier

\section{Structural Simulation Profiles}

The stock-price coordinate is represented by log-moneyness,
\begin{equation}
 m=\ln(S/S_0).
\end{equation}
Stock dependence can enter the closure through \(\Liq(S)\), \(\Aeff(S)\),
\(\nu(S)\), \(\gamma(S)\), and possibly the activity clock. The first
simulation holds \(\gamma\) constant and uses a simple maturity-dependent
activity profile,
\begin{equation}
 \alpha(T)
 =
 \alpha_0\left(1+a_{\mathrm{short}}\exp(-T/\tau_\alpha)\right).
 \label{eq:activity-profile}
\end{equation}
This keeps the source of \(m\)-dependence visible in each profile.

\begin{table}[ht]
\centering
\small
\caption{Structural profiles used to generate stock dependence in the simulated surfaces.}
\label{tab:profile-definitions}
\begin{tabular}{p{0.22\linewidth}p{0.26\linewidth}p{0.22\linewidth}p{0.20\linewidth}}
\toprule
Profile & Varies with \(m\) & Fixed across stock states & Purpose \\
\midrule
Liquidity-only & \(\Liq(m)\) &
\(\Aeff,z,\gamma,\alpha\) &
Isolates the inverse-square liquidity channel. \\
Resilience-response-only & \(z(m)=\nu(m)\Delta\) &
\(\Aeff,\Liq,\gamma,\alpha\) &
Isolates the response channel. \\
Combined structural & \(\Liq(m),z(m),\Aeff(m)\) &
\(\gamma\), simple \(\alpha(T)\) &
Combines stylised liquidity, resilience, and stress channels. \\
\bottomrule
\end{tabular}
\end{table}

The common functional forms are
\begin{align}
 \Liq(m)
 &=
 L_0\max\!\left(L_{\mathrm{floor}},
 1+L_{\mathrm{skew}}m+L_{\mathrm{smile}}m^2\right),
 \label{eq:liquidity-profile}\\
 z(m)
 &=
 z_{\mathrm{regime}}\exp(\nu_{\mathrm{skew}}m),
 \label{eq:z-profile}\\
 \Aeff(m)
 &=
 A_{\mathrm{eff},0}\exp(-A_{\mathrm{stress}}m).
 \label{eq:forcing-profile}
\end{align}
Positive \(L_{\mathrm{skew}}\) lowers liquidity for negative \(m\), while
positive \(L_{\mathrm{smile}}\) changes liquidity away from the money.
Positive \(\nu_{\mathrm{skew}}\) makes \(z(m)\) smaller on the downside.
In the combined profile, \(A_{\mathrm{stress}}\) raises the effective forcing
scale for negative \(m\). These choices are phenomenological simulation
controls, not fitted estimates.

\section{Numerical Implementation}

The 121 stock states range from $-0.30$ to $0.30$, and 96 maturities range from $1/12$ to 2 years. The surfaces are evaluated on a rectangular grid in \((m,T)\). For each grid
point, the selected structural profile determines \(\Liq(m)\), \(z(m)\), and
\(\Aeff(m)\); the selected response branch determines \(\Fresp_\gamma(z)\);
and Eqs.~\eqref{eq:operational-closure}--\eqref{eq:plotted-volatility}
determine the plotted volatility. The figures were generated with Python,
NumPy, and Matplotlib. Surface figures use boxed three-dimensional axes, and
the paired contour figures show black contour lines of the same surface.

\section{Relation To Calibrated Local Volatility}

The plotted object is best interpreted as a projected structural volatility
surface. It is local-volatility-like because it assigns a volatility level to
a state coordinate and maturity, but it is not a calibrated Dupire local
volatility surface. A calibrated local-volatility surface is usually inferred
from option prices under a risk-neutral pricing measure and must be consistent
with the absence of static arbitrage in the associated option surface. The
present construction instead starts from order-book structural quantities and
projects them into a surface morphology.

This distinction matters for interpretation. The simulation is designed to
ask whether plausible order-book channels can generate volatility-surface
shapes that resemble known market phenomenology, especially equity-like
downside skew. It does not claim that the resulting surface is directly
tradeable, arbitrage-free, or equivalent to a market-calibrated local
volatility input.

For comparison with option-calibrated surfaces, the state coordinate \(m\)
should be read as a log-moneyness coordinate. In a later empirical comparison,
one would normally align this with log-strike or log-forward moneyness, for
example \(m=\ln(K/S_0)\) or \(m=\ln(K/F_0(T))\), depending on the calibration
convention. The current notation \(m=\ln(S/S_0)\) keeps the connection to the
structural stock-price coordinate visible, while the plotted surface should
be compared only qualitatively with calibrated local volatility.

Within that qualitative interpretation, the combined structural high-\(z\)
case is the most market-like of the first simulations. It combines thinner
downside liquidity, weaker downside resilience through \(z(m)=\nu(m)\Delta\),
and stronger downside forcing. The low-\(z\) resilience-response-only case is
mainly diagnostic: at fixed \(\gamma\), the leading low-\(z\) approximation
does not depend on \(z\), so varying \(z(m)\) alone cannot generate leading
stock dependence in that branch.

\section{Parameter Choices}

The baseline values are dimensionless simulation controls. They set a
reasonable scale for visual comparison across profiles and regimes, but they
should not be read as empirical estimates.

\begin{table}[ht]
\centering
\small
\caption{Dimensionless baseline values and simulation ranges.}
\label{tab:parameter-ranges}
\begin{tabular}{p{0.18\linewidth}p{0.13\linewidth}p{0.13\linewidth}p{0.13\linewidth}p{0.30\linewidth}}
\toprule
Parameter & Baseline & Minimum & Maximum & Role \\
\midrule
\(\gamma\) & 0.55 & 0.30 & 0.80 & Long-memory exponent; avoids endpoints 0 and 1. \\
\(\Delta\) & 1.0 & 0.5 & 2.0 & Operational coarse-graining scale. \\
\(A_{\mathrm{eff},0}\) & 1.0 & 0.5 & 2.0 & Effective signed-forcing intensity scale. \\
\(L_0\) & 1.0 & 0.75 & 1.5 & Baseline liquidity slope. \\
\(z_{\mathrm{low}}\) & 0.05 & 0.01 & 0.10 & Low-\(z\) response-regime marker. \\
\(z_{\mathrm{high},0}\) & 10.0 & 5.0 & 25.0 & High-\(z\) response-regime marker. \\
\(\alpha_0\) & 1.0 & 0.5 & 2.0 & Baseline deterministic activity. \\
\(a_{\mathrm{short}}\) & 0.25 & 0.0 & 0.75 & Short-maturity activity lift. \\
\(\tau_\alpha\) & 0.50 & 0.25 & 1.0 & Activity decay scale in years. \\
\(L_{\mathrm{skew}}\) & 0.80 & 0.0 & 1.5 & Lower downside liquidity when positive. \\
\(L_{\mathrm{smile}}\) & 1.50 & 0.0 & 3.0 & Liquidity variation away from the center. \\
\(\nu_{\mathrm{skew}}\) & 1.00 & 0.0 & 2.0 & Weaker downside resilience when positive. \\
\(A_{\mathrm{stress}}\) & 0.75 & 0.0 & 1.5 & Stronger downside forcing in the combined profile. \\
\(L_{\mathrm{floor}}\) & 0.20 & 0.10 & 0.50 & Positive floor for the liquidity profile. \\
\bottomrule
\end{tabular}
\end{table}

\FloatBarrier

\section{Simulated Surfaces}

Each figure pair has the same interpretation. The left panel is the boxed
surface view of projected local volatility \(\sigloc(m,T)\). The right panel
is the black contour-line projection of the same surface. The horizontal
state coordinate is \(m=\ln(S/S_0)\), the maturity coordinate is \(T\) in
years, and the vertical surface coordinate is projected local volatility in
dimensionless simulation units.

\subsection{Liquidity-Only Profile}

\begin{figure}[ht]
 \centering
 \begin{subfigure}{0.48\linewidth}
  \centering
  \includegraphics[width=\linewidth]{\FigDir/mpl-surface3d-liquidity-only-low-z-v3.pdf}
  \caption{Boxed surface view.}
 \end{subfigure}
 \hfill
 \begin{subfigure}{0.48\linewidth}
  \centering
  \includegraphics[width=\linewidth]{\FigDir/mpl-contour-lines-liquidity-only-low-z-v3.pdf}
  \caption{Black contour-line projection.}
 \end{subfigure}
 \caption{Liquidity-only simulation in the low-\(z\) response regime. The only structural source of stock dependence is the liquidity profile \(\Liq(m)\), which enters the closure through the inverse-square factor \(1/\Liq(m)^2\). Since the low-\(z\) response approximation is independent of \(z\) at fixed \(\gamma\), the response channel does not contribute leading-order stock dependence in this panel.}
 \label{fig:liquidity-low-z}
\end{figure}

\begin{figure}[ht]
 \centering
 \begin{subfigure}{0.48\linewidth}
  \centering
  \includegraphics[width=\linewidth]{\FigDir/mpl-surface3d-liquidity-only-high-z-v3.pdf}
  \caption{Boxed surface view.}
 \end{subfigure}
 \hfill
 \begin{subfigure}{0.48\linewidth}
  \centering
  \includegraphics[width=\linewidth]{\FigDir/mpl-contour-lines-liquidity-only-high-z-v3.pdf}
  \caption{Black contour-line projection.}
 \end{subfigure}
 \caption{Liquidity-only simulation in the high-\(z\) response regime. The liquidity profile remains the only source of stock dependence because \(z\) is held constant in this profile. The comparison with Fig.~\ref{fig:liquidity-low-z} isolates the change in response-regime scale while preserving the same liquidity-driven surface shape.}
 \label{fig:liquidity-high-z}
\end{figure}

\FloatBarrier

\subsection{Resilience-Response-Only Profile}

\begin{figure}[ht]
 \centering
 \begin{subfigure}{0.48\linewidth}
  \centering
  \includegraphics[width=\linewidth]{\FigDir/mpl-surface3d-resilience-response-only-low-z-v3.pdf}
  \caption{Boxed surface view.}
 \end{subfigure}
 \hfill
 \begin{subfigure}{0.48\linewidth}
  \centering
  \includegraphics[width=\linewidth]{\FigDir/mpl-contour-lines-resilience-response-only-low-z-v3.pdf}
  \caption{Black contour-line projection.}
 \end{subfigure}
 \caption{Resilience-response-only simulation in the low-\(z\) response regime. The profile varies \(z(m)=\nu(m)\Delta\), but the leading low-\(z\) response approximation in Eq.~\eqref{eq:low-z-response} does not depend on \(z\) at fixed \(\gamma\). The figure is therefore a diagnostic check: it shows the maturity activity effect without an additional leading-order response-driven stock dependence.}
 \label{fig:resilience-low-z}
\end{figure}

\begin{figure}[ht]
 \centering
 \begin{subfigure}{0.48\linewidth}
  \centering
  \includegraphics[width=\linewidth]{\FigDir/mpl-surface3d-resilience-response-only-high-z-v3.pdf}
  \caption{Boxed surface view.}
 \end{subfigure}
 \hfill
 \begin{subfigure}{0.48\linewidth}
  \centering
  \includegraphics[width=\linewidth]{\FigDir/mpl-contour-lines-resilience-response-only-high-z-v3.pdf}
  \caption{Black contour-line projection.}
 \end{subfigure}
 \caption{Resilience-response-only simulation in the high-\(z\) response regime. Here \(z(m)\) enters directly through \(K_\gamma^{(0)}z^{\gamma-1}\). With \(\gamma<1\), lower downside \(z(m)\) increases the response multiplier, producing stock dependence even though liquidity and forcing are held fixed.}
 \label{fig:resilience-high-z}
\end{figure}

\FloatBarrier

\subsection{Combined Structural Profile}

\begin{figure}[ht]
 \centering
 \begin{subfigure}{0.48\linewidth}
  \centering
  \includegraphics[width=\linewidth]{\FigDir/mpl-surface3d-combined-structural-low-z-v3.pdf}
  \caption{Boxed surface view.}
 \end{subfigure}
 \hfill
 \begin{subfigure}{0.48\linewidth}
  \centering
  \includegraphics[width=\linewidth]{\FigDir/mpl-contour-lines-combined-structural-low-z-v3.pdf}
  \caption{Black contour-line projection.}
 \end{subfigure}
 \caption{Combined structural simulation in the low-\(z\) response regime. The surface combines liquidity variation and downside forcing while the leading response approximation remains independent of \(z\) at fixed \(\gamma\). The resulting shape should be interpreted as a stylised liquidity-and-forcing effect rather than as a calibrated volatility surface.}
 \label{fig:combined-low-z}
\end{figure}

\begin{figure}[ht]
 \centering
 \begin{subfigure}{0.48\linewidth}
  \centering
  \includegraphics[width=\linewidth]{\FigDir/mpl-surface3d-combined-structural-high-z-v3.pdf}
  \caption{Boxed surface view.}
 \end{subfigure}
 \hfill
 \begin{subfigure}{0.48\linewidth}
  \centering
  \includegraphics[width=\linewidth]{\FigDir/mpl-contour-lines-combined-structural-high-z-v3.pdf}
  \caption{Black contour-line projection.}
 \end{subfigure}
 \caption{Combined structural simulation in the high-\(z\) response regime. This profile combines the inverse-square liquidity channel, the response multiplier \(K_\gamma^{(0)}z^{\gamma-1}\), and a stylised downside forcing profile. The figure is the closest of the six simulations to the intended phenomenological comparison with a calibrated volatility surface, while remaining a mechanism simulation rather than an empirical fit.}
 \label{fig:combined-high-z}
\end{figure}

\FloatBarrier

\section{Robustness Checks}

The numerical checks verify that the simulated surfaces are finite and
positive on the grid, that the liquidity floor prevents invalid liquidity
values, and that the response-regime cases behave consistently with their
asymptotic formulas. The checks also verify the qualitative distinction
between the low-\(z\) and high-\(z\) response channels: when only \(z(m)\)
varies, the low-\(z\) branch has no leading-order \(m\)-dependence at fixed
\(\gamma\), while the high-\(z\) branch does.

\begin{figure}[ht]
\centering
\includegraphics[width=.70\linewidth]{\FigDir/v2/response-crossover.pdf}
\caption{\textbf{Frozen variance crossover.} The scaled increment variance and running rate slope use \(\gamma=1/2\). At fixed finite cutoff the microscopic, intermediate and saturated regimes have different lag dependence. Each plotted response uses its stated cutoff.}
\end{figure}
\FloatBarrier

For the resilience-only profile, the maximum stock-state spread at fixed
maturity is zero in the low-\(z\) branch and 0.11837 in the high-\(z\)
branch. In the combined high-\(z\) case the corresponding spread is
0.63977. These are spreads of the plotted volatility on the finite grid.
In the sensitivity sweep over \(\Delta\), the code holds the assigned
\(z(m)\) profile fixed. The sweep therefore tests the explicit
\(\Delta^{-\gamma}\) factor at fixed \(z\); a sweep at fixed \(\nu\)
would also move along the response curve.

\section{A translating liquidity profile}
Write \(x\) for log price and let \(y(u)\) be the observed zero of the
imbalance field. For \(D,\nu,\lambda_{\rm dep}>0\), take
\begin{equation}
 \partial_u\Phi=D\Phi_{xx}-\nu\Phi-
 \lambda_{\rm dep}\operatorname{sgn}(x-y)+b(u)h(x-y),\qquad \Phi(y,u)=0 .
 \label{eq:auxiliary-field}
\end{equation}
Set \(a=\sqrt{\nu/D}\) and \(\Liq=\lambda_{\rm dep}/\sqrt{D\nu}\).
In the absence of forcing there is a family of steady profiles
\[
 \Phi_*(\xi)=-\frac{\lambda_{\rm dep}}{\nu}\operatorname{sgn}(\xi)
       (1-e^{-a|\xi|}),\qquad
 \psi(\xi)=-\Phi_*'(\xi)=\Liq e^{-a|\xi|},\qquad \xi=x-q .
\]
The observed-root variation generated by a field perturbation \(f\) is
\(p=f(0)/\Liq\). Differentiating the sign source therefore adds
\(2\lambda_{\rm dep}f(0)\delta_0/\Liq\) to the pinned operator.
An even bound mode \(e^{-k|x|}\) has derivative jump \(-2k\).
Continuity and the jump condition give \(Dk=\lambda_{\rm dep}/\Liq\),
so \(k=a\) and its eigenvalue is \(Dk^2-\nu=0\). There is only one
decaying solution of that jump equation. The free half-line modes give
the shape continuum \(( -\infty,-\nu]\). The gap concerns the transverse
spatial spectrum. The time spectrum of a driven root also depends on its forcing.

The spectral phase \(\langle\psi,f\rangle/\|\psi\|^2\) and \(p\) are
distinct observables. For \(f=e^{-2a|x|}-(2/3)e^{-a|x|}\) the phase
vanishes while \(p=1/(3\Liq)\). For aligned placement
\(h_*=\kappa_*\psi\), \(\kappa_*=a/(2\Liq)\), substitution in
Eq.~\eqref{eq:auxiliary-field} instead gives an exact invariant family
\(\Phi(x,u)=\Phi_*(x-q(u))\), with \(y=q\) and \(\dot q=\kappa_* b\)
for initial data on that family. At unit parameters the finest cross-root
weak residual was \(3.16\times10^{-5}\); shifting the placement by 0.4
raised it to 0.218 for the same trial profile. The shifted-profile check
tests the equation for a translate, rather than measuring the trajectory
of a perturbed solution.

\begin{algorithm}[htbp]
\caption{Recovering the observed zero and testing the translating field}
\small
\textbf{Input:} Increasing spatial grid, translated profile samples at three operational times, spatial forcing profile and test function.\\
\textbf{Output:} Observed zero, local slope and weak field residual.
\begin{enumerate}\setlength{\itemsep}{0.1em}
\item Count all crossings. Require one interior positive-to-negative crossing and reject ambiguous or reversed crossings.
\item Recover the zero by linear interpolation and estimate the slope from the bracketing samples. Record the error against the imposed translation separately.
\item Split the spatial quadrature at the recovered zero. Difference the field in time and integrate the diffusion term against the second derivative of the compact test function.
\item Repeat after displacing the forcing profile at fixed trial trajectory. Compare the two residuals at the same space and time steps.
\end{enumerate}
\end{algorithm}
\FloatBarrier
\begin{figure}[ht]
\centering
\includegraphics[width=.91\linewidth]{\FigDir/v2/translation-modes.pdf}
\caption{\textbf{Translation tangent and moving zero.} The normalised principal lattice mode is compared with the analytic tangent (left). On the invariant family the imposed forcing displaces the full profile and its zero together (right). The independent cross-root residual tests this identity in the field equation.}
\end{figure}
\FloatBarrier

\subsection{Implementation and resolution}
The finite-difference operator uses a central defect of strength
\(2\lambda_{\rm dep}/(\Liq\,\delta x)\), with Dirichlet walls one cell
beyond the outermost grid points. At \(D=\nu=\lambda_{\rm dep}=1\)
the largest eigenvalues on successively refined grids with
\(\delta x=0.16,0.08,0.04\) are \(-0.00631937,-0.00159490,-0.00039968\).
The infinite-lattice prediction is
\[
 -\nu\frac{(a\,\delta x)^2}{(1+\sqrt{1+(a\,\delta x)^2})^2}.
\]
The shape threshold and the translation eigenvalue are also tested on
boxes of half-width 6, 12 and 24. For the aligned weak-field comparison
we use three compact test functions and four operational times.
Refining \((\delta x,\delta u)\) from \((0.08,0.04)\) to \((0.01,0.005)\)
reduces the maximum weak residual from 0.00249 to 0.0000316. On the
fine grid the observed-root interpolation error is below 0.0000048.

\section{Causal liquidity reassessment}
Let \(v=\dot y\) on the aligned family and prescribe
\[
 \dot z_r=-rz_r+v,\qquad v=\kappa_*m_u-g\int_0^\infty \rho_\theta(r)z_r\,dr,
 \qquad \rho_\theta(r)=\frac{\sin(\pi\theta)}{\pi}r^{-\theta},
 \quad 0<\theta<\tfrac12 .
\]
Here \(g\ge0\) is the feedback strength. The causal response is measured
from zero memory. The Euler integral gives
\(K(s)=\int\rho_\theta(r)/(s+r)\,dr=s^{-\theta}\) for \(\Re s>0\).
Consequently the velocity transfer and the root impulse response obey
\begin{equation}
 H(s)=\frac{\kappa_*}{1+gK(s)},\qquad
 \widehat G(s)=\frac{H(s)}{s},\qquad
 G(u)=\kappa_* E_\theta(-gu^\theta),\quad u\ge0 .
 \label{eq:feedback}
\end{equation}
The factor \(1/s\) integrates the velocity once. Thus velocity under a
unit step also measures the root response to a unit impulse. We evaluate
the continuum function independently from its positive relaxation density;
see \href{https://arxiv.org/abs/1305.0161}{arXiv:1305.0161}.

For rates between \(r_-\) and \(r_+\), take midpoint nodes in \(\log r\)
with weights \(w_j=\rho_\theta(r_j)r_j\,\delta\log r\).
Replacing \(K\) by \(K_N(s)=\sum_jw_j/(s+r_j)\) gives
\(H_N(0)=\kappa_* /(1+g\sum_j w_j/r_j)>0\) for \(\kappa_*>0\).
The finite model therefore retains a zero-frequency response. The matrix
of decay rates is similar to \(\operatorname{diag}(r_j)+g dd^\top\),
where \(d_j=\sqrt{w_j}\); all its eigenvalues are positive.
The stability check and the value at zero frequency concern different
properties of the finite system.

For \(D=\nu=\lambda_{\rm dep}=g=1\), \(\kappa_*=1/2\). We used
\(\theta\in\{1/4,1/3,2/5\}\), independently varied \(r_-\in\{10^{-4},10^{-7}\}\),
\(r_+\in\{10^3,10^6\}\), and used 48 or 96 nodes and 17 logarithmic
frequencies from 0.01 to 1. At 96 nodes the maximum relative error against
adaptive quadrature at the same endpoints was \(2.46\times10^{-5}\).
Even on the widest rate interval the discrepancy from the infinite-rate
continuum reached 1.48\%. For the time comparison we used 64 nodes from
\(10^{-6}\) to \(10^5\), \(\theta=1/3\) and 25 logarithmic lags from
0.1 to 10. Implicit time integration and independent finite-pole inversion
differed by at most \(4.02\times10^{-9}\) in absolute response. The
finite-rate response differed from the continuum by up to 0.00370.

\subsection{Eliminating fast reassessment rates}
For a cutoff \(\Lambda\), retain all rates below it. Integrating the fast
states with zero initial history gives the exact causal memory
\(k_>(u)=\sum_{r_j\ge\Lambda} w_j e^{-r_ju}\) acting on the past velocity.
For \(\Re s\ge0\) the expansion and its remainder obey
\[
 K_>(s)=M_0-sM_1+R(s),\quad
 M_n=\sum_{r_j\ge\Lambda}w_j/r_j^{n+1},\quad |R(s)|\le |s|^2 M_2 .
\]
The estimate follows by expanding each resolvent with exact remainder
\(s^2/[r_j^2(r_j+s)]\). The reduced denominator is
\(1+g[K_<(s)+M_0-sM_1]\), so the relative transfer error is bounded by
\(g|s|^2 M_2/|1+g(K_<+M_0-sM_1)|\) wherever the denominator is nonzero.
The unit of operational time and the physical values of \(g,\kappa_*\)
are fixed; the expansion parameter is \(|s|/\Lambda\).
The derivative approximation applies only on its retained band.

\begin{table}[ht]
\centering\small
\caption{Maximum relative velocity-transfer error on 25 logarithmic frequencies
from 0.001 to 0.2. The three approximations use the same 64 rates and gain.}
\begin{tabular}{rrrr}
\toprule
\(\Lambda\) & Omit shell & Retain \(M_0\) & Retain \(M_0-sM_1\) \\
\midrule
2  & 0.2990 & 0.005200 & 0.0002532 \\
5  & 0.2135 & 0.001811 & 0.0000400 \\
20 & 0.1321 & 0.000372 & 0.00000251 \\
\bottomrule
\end{tabular}
\end{table}
\begin{algorithm}[htbp]
\caption{Causal root response and fast-rate elimination}
\small
\textbf{Input:} Rate endpoints, number of nodes, \(\theta,g,\kappa_*\), cutoff \(\Lambda\), operational lags and frequencies.\\
\textbf{Output:} \(G(u)\), complex \(H(i\omega)\), and errors for each shell approximation.
\begin{enumerate}\setlength{\itemsep}{0.1em}
\item Construct positive midpoint weights in log rate. Evaluate the complex transfer from the weighted resolvent and compare it with adaptive quadrature over the same endpoints.
\item Starting from zero memory, integrate the mode equations under a unit step using implicit Radau time stepping. Reconstruct the same response independently from the symmetric decay matrix. Use scalar quadrature of the continuum relaxation spectrum as a third comparison.
\item Split at \(r_j\ge\Lambda\). Accumulate \(M_0,M_1,M_2\) only over the fast modes. Keep the slow memory exact and use the original physical units.
\item At each retained frequency compare the full transfer with the exact shell, omission, static memory and first derivative. Record the complex relative error and its analytic bound.
\end{enumerate}
\end{algorithm}
\FloatBarrier
\begin{figure}[htbp]
\centering
\includegraphics[width=.9\linewidth]{\FigDir/v2/liquidity-feedback.pdf}
\caption{Left: the root displacement after a unit impulse in signed input, equivalently the velocity under a unit step, at 25 operational lags. The circles use 64 finite rates at \(\theta=1/3\) and the line is the continuum relaxation. Right: relative error in the complex velocity transfer after eliminating rates \(r_j\ge5\), over the retained frequency band. The static and first-derivative approximations use the same fast modes as the omission control.}
\label{fig:feedback}
\end{figure}
Restoring the exact shell memory reproduces the full transfer to within
\(4.50\times10^{-16}\) on the test grid. For input spectrum
\(S_m(\omega)\sim A|\omega|^{\gamma-1}\) the ideal continuum predicts
\(S_v(\omega)\sim A(\kappa_*/g)^2|\omega|^{2\theta+\gamma-1}\) for
\(g>0\). Setting \(2\theta+\gamma=1\) is a choice of the prescribed
rate density. Testing an intermediate diffusive regime requires an
independent estimate of forcing memory, the slope of the integrated
response and the variance of the observed root over the same lags. A
finite rate floor changes the asymptote. A physical selection principle
for the feedback spectrum and transverse stability of a stochastic moving
front remain open questions.

\section{Response balance on an operational-lag window}
\label{sec:response-balance}

Let \(m_u\) denote the signed input, \(v\) the velocity of the aligned zero and
\(G(u)\) its displacement following a unit impulse in that input. Integrating
the velocity gives \(\widehat G(s)=H(s)/s\). Under zero prehistory this is
also the transform of the velocity following a unit step. For a stationary
input with \(S_m(\omega)\propto |\omega|^{\gamma-1}\), linear filtering
gives \(S_v=|H(i\omega)|^2 S_m\). In the ideal continuum, if
\(G(u)\sim c u^{-\beta}\) with \(0<\beta<1\) and regular variation,
then \(H(s)\sim c\Gamma(1-\beta)s^\beta\) on the positive Laplace axis.
For the causal continuum law derived above its explicit transform gives the
same power on the frequency axis. The corresponding velocity-spectrum
exponent is \(2\beta+\gamma-1\). On a finite window we measure
\[
 \beta_{\rm eff}(\Delta)=-\frac{d\log G(\Delta)}{d\log\Delta},
 \qquad B(\Delta)=2\beta_{\rm eff}(\Delta)+\gamma-1 .
\]
These local measurements compare two separately observed slopes. Even a
small residual over a decade does not determine the increment variance of
the observed zero. With a finite slow-rate floor the velocity transfer
ultimately approaches a positive constant.

\subsection{Forcing paths and the response measurement}
On a periodic grid of duration \(T_u=262144\), with operational spacing
\(h=1\) or \(1/2\), the positive-frequency coefficients of real Gaussian
white noise are multiplied by
\( |2\sin(\omega h/2)/h|^{(\gamma-1)/2}\).
The zero coefficient is set to zero. A real inverse Fourier transform
produces a stationary periodic sample with its sample mean fixed. There is
no start-up transient in this construction. On each of 32 independent
paths per exponent and spacing we regress its realised log periodogram on
\(\log\omega\), add one to the fitted slope, and repeat using every second
bin of the same band. The angular frequencies are \([1/320,1/32]\):
the convention here is reciprocal lag, rather than cycles per lag. On the
finer grid using all bins, the mean estimates are 0.39762 and 0.62319 for
prescribed inputs 0.4 and 0.6. Independent paths determine sampling
uncertainty; Fourier bins from a single path do not count as replicates.

The response experiment uses the finite transfer and its causal inverse
rather than fitting a price-variance slope. The primary parameters are
\(g=4,\ \kappa_*=1/2\), rate endpoints \(10^{-7},100\), and 128
logarithmic midpoint cells. For each input law, the prescribed rate
exponents are \((1-\gamma)/2\) and that value shifted by \(\pm0.08\).
At 17 log-spaced centres from 32 to 320 we regress \(\log G\) on log lag
over \([\Delta/\sqrt2,\sqrt2\Delta]\). Thus the centre interval and the
extended support used by its endpoint regressions are distinct. Samples
must be positive over the entire regression support.

\subsection{A change in the placement of signed forcing}
For the linearised shape comparison use \(D=\nu=\mathcal L=1\).
Write \(k=\sqrt{1+s}\), taking the branch with positive real part.
For \(h_\epsilon(x)=[(1-\epsilon)e^{-|x|}+\epsilon e^{-2|x|}]/2\)
the forcing at the origin stays fixed. The point-interaction resolvent at
the origin is \(\int e^{-k|x|}h_\epsilon(x)\,dx/[2(k-1)]\).
The root first variation divides the field perturbation by
\(\mathcal L\); multiplication by \(s=(k-1)(k+1)\) gives the velocity
mobility
\[
 M_\epsilon(s)=\kappa_* A_\epsilon(s),\quad
 A_\epsilon(s)=1-\epsilon/(k+2),\qquad
 H_\epsilon(s)=\frac{\kappa_* A_\epsilon(s)}
 {1+g A_\epsilon(s)K_N(s)} .
\]
Here the feedback gain is held fixed as the profile changes, and the factor
\(A_\epsilon\) multiplies the feedback as well as the external forcing.
For \(\epsilon=0\) this recovers the aligned model. For
\(\epsilon=0.2\), \(A_\epsilon\) has positive real part on
\(\Re s\ge0\), and its imaginary part has the sign of \(\Im s\).
The positive Stieltjes sum \(K_N\) has positive real part and the opposite
imaginary sign. Consequently \(\Re(A_\epsilon K_N)>0\), so the feedback
denominator cannot vanish there. This placement control concerns the
first variation of the observed zero. Its extension to a nonlinear
random front requires separate root tracking.

\begin{algorithm}[htbp]
\caption{Independent forcing-memory and integrated-response estimates}
\small\textbf{Input:} Operational lag centres, forcing realisations, feedback parameters, rate bounds and forcing placement.\\
\textbf{Output:} Input-memory intervals, response slopes, balance intervals and whole-window classifications.
\begin{enumerate}\setlength{\itemsep}{0.1em}
\item Fit each path at both Fourier resolutions. Form Student intervals across 32 paths at each spacing, resolution and input exponent. Use the two-sided critical value \(t_{31,1-0.05/16}\) for eight comparisons. For each input exponent keep the envelope of its four intervals.
\item Evaluate \(\widehat G=H_\epsilon/s\) on the Bromwich line at \(s=(A/2+i\pi j)/u\). Multiply alternating real parts by \(e^{A/2}/u\), halve the zero term and Euler-average partial sums numbered \(n\) to \(n+m\) using weights \(2^{-m}\binom{m}{j}\) indexed from zero. Use \((A,n,m)=(24,40,30)\) and \((28,52,36)\).
\item Repeat with 64 and 128 rate cells and 9 and 17 log-lag samples in each fit. Compare aligned inversions with finite pole inversion. Keep the minimum and maximum response slope across refinements at each centre.
\item Combine endpoints monotonically, \([2\beta_{\min}+\gamma_{\rm lo}-1,\,2\beta_{\max}+\gamma_{\rm hi}-1]\). Compatibility requires inclusion in \([-0.1,0.1]\) at every centre. An offset is resolved only when all intervals lie strictly beyond the same edge. Retain other results as inconclusive and flag nonpositive responses.
\end{enumerate}
\end{algorithm}
The input intervals describe finite-window estimates under the periodic
forcing model. The response envelope measures observed numerical
sensitivity. It supplies no coverage guarantee for continuum truncation
or misspecification. At all tested aligned settings the largest absolute
difference from independent finite pole inversion was
\(3.32\times10^{-10}\) or less.

\begin{table}[htbp]
\centering\small
\caption{Primary response experiment. The interval runs from the smallest lower endpoint to the largest upper endpoint over all 17 lag centres. Labels use the predeclared whole-window criterion.}
\begin{tabular}{rrrrl}
\toprule
Input \(\gamma\) & Rate \(\theta\) & Min. lower & Max. upper & Outcome\\
\midrule
0.4 & 0.22 & $-0.219$ & $-0.116$ & Negative offset\\
0.4 & 0.30 & $-0.071$ & $ 0.028$ & Compatible\\
0.4 & 0.38 & $ 0.084$ & $ 0.185$ & Inconclusive\\
0.6 & 0.12 & $-0.185$ & $-0.047$ & Inconclusive\\
0.6 & 0.20 & $-0.051$ & $ 0.080$ & Compatible\\
0.6 & 0.28 & $ 0.095$ & $ 0.219$ & Inconclusive\\
\bottomrule
\end{tabular}
\end{table}
\begin{figure}[htbp]
\centering\includegraphics[width=.93\linewidth]{\FigDir/v2/response-balance.pdf}
\caption{Measured balance of the aligned response for the two input laws. Curves give the finest-grid estimates for each prescribed rate exponent. Coloured bands combine the forcing intervals with observed response refinements. The grey strip marks the fixed compatibility tolerance. A signed estimate can be distinguishable from zero while its interval still overlaps the tolerance edge.}
\label{fig:response-balance}
\end{figure}
\FloatBarrier
Both primary matched cases satisfy the tolerance. Of the four primary
offset controls, only the lower offset at input exponent 0.4 clears the
separation threshold throughout the decade. At each matched setting
raising the slow-rate floor to 0.02 produces a negative offset, as does
turning feedback off. Reducing the feedback strength from four to one
gives an inconclusive result for both matches. Raising the rate ceiling
to 1000 or changing the linearised forcing profile with \(\epsilon=0.2\)
preserves compatibility for both. The full sweep includes 36
configurations and 612 lag-centre results. This demonstrates finite-window
compatibility in a prescribed construction. It leaves the independent
actual-root variance test on the same lag grid open.

\section{Limitations And Interpretation}

The simulations are structural illustrations. They do not estimate
\(\Liq(S)\), \(\nu(S)\), \(\Aeff(S)\), \(\gamma(S)\), or \(\alpha(T)\) from
market data. The stock dependence of the plotted surfaces is selected through
declared phenomenological profiles: inverse-square liquidity effects,
response dependence through \(z(S)=\nu(S)\Delta\), and a stylised downside
forcing scale. The role of the figures is therefore to make the theory
visually testable and to clarify which structural channels can plausibly
produce surface morphology before any empirical calibration is attempted.

\section*{Use of Generative AI}

Generative AI tools were used as computational and editorial assistance in
preparing this supplementary note, including drafting explanatory text,
organizing simulation code, and checking figure and table presentation. The
theoretical interpretation, modelling assumptions, parameter choices, code
outputs, and final scientific claims in the historical edition were reviewed by the author(s). The additional calculations in this working edition are presented for scientific review. No empirical market data are used in
these simulations; the figures are generated from the stated equations and
declared parameter profiles.

\end{document}
